\section{Weight Transfer \& Load Distribution}

This section documents the calculation of normal loads on each wheel as a function of vehicle acceleration and cornering forces. Weight transfer is fundamental to the tire model and affects grip levels. The implementation has been verified against the source code and cross-checked for mathematical correctness.

\cite{milliken1995race} provides the classical treatment of weight transfer in vehicle dynamics. The implementation here is grounded in Newton's second law applied in a rotating reference frame, simplified for real-time computation.

\subsection{Total Vehicle Weight}

\begin{equation}
  W = m \cdot g
  \label{eq:total-weight}
\end{equation}

\textbf{Physical Meaning}: The total weight of the vehicle is the product of mass and gravitational acceleration. This is split equally between the front and rear axles in a 50/50 distribution (assumed balanced CoG).

\textbf{Parameters}:
\begin{itemize}
  \item $m = \mathrm{carMassKg}$: Vehicle mass in kg (default: 1500 kg)
  \item $g = \mathrm{GRAVITY}$: Gravitational acceleration (9.8 m/s²)
\end{itemize}

\subsection{Longitudinal Weight Transfer}

\begin{equation}
  \Delta W_{\text{long}} = m \cdot a_{\text{long}} \cdot \frac{h_{\text{cog}}}{L_{\text{wb}}}
  \label{eq:longitudinal-transfer}
\end{equation}

\textbf{Physical Meaning}: Longitudinal acceleration (braking/acceleration) causes weight to transfer between the front and rear axles. The transfer is proportional to acceleration magnitude, center of gravity height, and inverse wheelbase.

\textbf{Source Code}: \coderef{weight.js}{25--35}

\begin{lstlisting}[caption={Longitudinal weight transfer computation}]
export function computeWeightTransfer(dt) {
  const body = state.body;
  const params = state.params;
  const mass = params.carMassKg;
  const cogHeight = params.cogHeight;
  const wheelbase = params.wheelbase;
  const trackWidth = params.trackWidth;

  // Clamp acceleration to ±2g to prevent constraint solver transients
  const clampedLongAccel = clamp(body.accelX * forwardUnitX + body.accelY * forwardUnitY,
                                  -2 * GRAVITY, 2 * GRAVITY);
  const longitudinalTransfer = mass * clampedLongAccel * cogHeight / wheelbase;
\end{lstlisting}

\textbf{Notes}:
\begin{itemize}
  \item Acceleration is clamped to $\pm 2g$ to prevent unrealistic weight transfer from constraint solver numerical artifacts
  \item CoG height typically 0.3–0.5 m above ground
  \item Wheelbase typically 2.4–2.8 m for passenger cars
\end{itemize}

\subsection{Lateral Weight Transfer}

\begin{equation}
  \Delta W_{\text{lat}} = m \cdot a_{\text{lat}} \cdot \frac{h_{\text{cog}}}{T_{\text{w}}}
  \label{eq:lateral-transfer}
\end{equation}

\textbf{Physical Meaning}: Lateral acceleration (cornering) causes weight to transfer from inside to outside wheels, proportional to track width. Wider track width reduces weight transfer for the same lateral acceleration.

\textbf{Source Code}: \coderef{weight.js}{36--40}

\begin{lstlisting}[caption={Lateral weight transfer computation}]
  const lateralAccel = body.accelX * rightUnitX + body.accelY * rightUnitY;
  const lateralTransfer = mass * lateralAccel * cogHeight / trackWidth;

  state.wheelLoads.longitudinalTransfer = longitudinalTransfer;
  state.wheelLoads.lateralTransfer = lateralTransfer;
\end{lstlisting}

\subsection{Per-Wheel Normal Loads}

\begin{align}
  N_{\text{FL}} &= \max\left(0, \frac{W}{4} - \Delta W_{\text{long}} - \Delta W_{\text{lat}}\right) \\
  N_{\text{FR}} &= \max\left(0, \frac{W}{4} - \Delta W_{\text{long}} + \Delta W_{\text{lat}}\right) \\
  N_{\text{RL}} &= \max\left(0, \frac{W}{4} + \Delta W_{\text{long}} - \Delta W_{\text{lat}}\right) \\
  N_{\text{RR}} &= \max\left(0, \frac{W}{4} + \Delta W_{\text{long}} + \Delta W_{\text{lat}}\right)
  \label{eq:wheel-loads}
\end{align}

\textbf{Physical Meaning}: Each wheel load is computed from the baseline quarter-weight ($W/4$), adjusted by transfer amounts. Negative loads are clamped to zero, representing wheel liftoff.

\textbf{Source Code}: \coderef{weight.js}{42--55}

\begin{lstlisting}[caption={Per-wheel normal load computation}]
const halfWeight = (mass * GRAVITY) / 2;
const loads = {
  frontLeft:  Math.max(0, halfWeight * 0.5 - longitudinalTransfer - lateralTransfer),
  frontRight: Math.max(0, halfWeight * 0.5 - longitudinalTransfer + lateralTransfer),
  rearLeft:   Math.max(0, halfWeight * 0.5 + longitudinalTransfer - lateralTransfer),
  rearRight:  Math.max(0, halfWeight * 0.5 + longitudinalTransfer + lateralTransfer),
};
\end{lstlisting}

\textbf{Notes}:
\begin{itemize}
  \item These loads directly scale Pacejka tire model output (Section \ref{sec:pacejka-forces})
  \item Wheel liftoff occurs when $\Delta W_{\text{long}}$ or $\Delta W_{\text{lat}}$ exceed quarter-weight
  \item Weight transfer is critical for accurate grip simulation
  \item The clamping to zero prevents negative normal loads, which would be physically impossible
\end{itemize}

\newpage
